\section*{Hoofdstuk \ref{ch:g6:life-through-seasons} --- Leven door de seizoenen}

\begin{solution}{exo:g6:life-through-seasons:1}
De daglengte (het licht), de \omterm{def:g6:exploring-our-environment:conditions}{temperatuur}, en de \omterm{def:g6:exploring-our-environment:conditions}{vochtigheid} van bodem of
lucht in het spoor van regen en droogte.
\end{solution}

\begin{solution}{exo:g6:life-through-seasons:2}
\omterm{def:g6:life-through-seasons:trees}{Bladverliezend} --- laat in de herfst al zijn \omterm{def:g1:plants-around-us:parts}{bladeren} vallen: eik, beuk,
paardenkastanje. \omterm{def:g6:life-through-seasons:trees}{Groenblijvend} --- houdt de hele winter werkende
\omterm{def:g1:plants-around-us:parts}{bladeren}: den, hulst.
\end{solution}

\begin{solution}{exo:g6:life-through-seasons:3}
Een piepkleine complete scheut --- miniatuurbladeren opgevouwen in
weerbestendige \omterm{def:g1:animals-around-us:covering}{schubben}. Hij is de zomer ervoor gebouwd.
\end{solution}

\begin{solution}{exo:g6:life-through-seasons:4}
De eenjarige gaat helemaal dood en brengt de winter door als \omterm{def:g2:from-seed-to-plant:seed}{zaad}. Bij de
narcis sterven de bovengrondse delen af terwijl de \omterm{def:g1:plants-around-us:plant}{plant} ondergronds als
\omterm{ex:g4:plants-without-seeds:bulbtuber}{bol} wacht.
\end{solution}

\begin{solution}{exo:g6:life-through-seasons:5}
Omdat de scheuten in zijn \omterm{def:g6:life-through-seasons:buds}{knoppen} vorige zomer al af waren; de twijg
miste alleen warmte, en die levert het huis weken eerder dan de tuin.
\end{solution}

\begin{solution}{exo:g6:life-through-seasons:6}
\omterm{def:g3:animals-through-seasons:migration}{Trekken}: de zwaluw. \omterm{def:g3:animals-through-seasons:hibernation}{Winterslaap}: de egel (of de afgesloten slak). Actief
blijven: de vos, de ekster, het roodborstje. Overwinteren als jong
stadium --- \omterm{def:g2:animals-grow-up:born}{ei}, larve of \omterm{ex:g2:animals-grow-up:butterfly}{pop}: de meeste \omterm{prop:g3:sorting-living-things:groups}{insecten}, zoals de sprinkhaan of
de vlinder.
\end{solution}

\begin{solution}{exo:g6:life-through-seasons:7}
Vlinder: een \omterm{ex:g2:animals-grow-up:butterfly}{pop} onder een afdeksteen. Sprinkhaan: eitjes in de bodem.
Lieveheersbeestje: een slapende volwassene, diep weggedrukt in de
bladlaag. Slak: afgesloten in haar huisje in de berm. Zwaluw: in Afrika,
als zichzelf, getrokken.
\end{solution}

\begin{solution}{exo:g6:life-through-seasons:8}
Zodat elk verschil tussen de vier tabellen van de seizoenen komt en niet
van het onderzoeken --- de regel van de eerlijke proef uit
\cref{met:g6:exploring-our-environment:survey}, toegepast over de tijd in
plaats van over locaties.
\end{solution}

\begin{solution}{exo:g6:life-through-seasons:9}
Ten eerste zijn veel ``afwezigen'' aanwezig in verborgen vormen ---
eitjes, poppen, slapende volwassenen, bollen --- die een onderzoek dat
alleen rondloopt mist zonder graven en optillen. Ten tweede zijn sommige
niet verloren maar elders: de trekkers komen terug. De januarilijst geeft
de bewoners te laag weer en telt reizigers als verdwenen.
\end{solution}

\begin{solution}{exo:g6:life-through-seasons:10}
De \omterm{def:g1:plants-around-us:plant}{plant} zit in haar ondergrondse opslagtoestand --- levend, wachtend als
bollen. In maart sturen de bollen weer \omterm{def:g1:plants-around-us:parts}{bladeren} en \omterm{def:g3:flowers-fruits-seeds:parts}{bloemen} omhoog,
dichter als de pollen zich vermeerderd hebben.
\end{solution}

\begin{solution}{exo:g6:life-through-seasons:11}
Weg waarheen, als wat? Kikkers overwinteren in rust --- ingegraven in de
modderbodem van de vijver of in vochtige schuilplaatsen op het land ---
en brengen de koude maanden verstijfd door tot de lente ze terugroept
naar het water om zich voort te \omterm{def:g1:plants-around-us:plant}{planten}.
\end{solution}

\begin{solution}{exo:g6:life-through-seasons:12}
Het winterprobleem van een werkend \omterm{def:g1:plants-around-us:parts}{blad} is water vasthouden: bevroren
grond geeft de \omterm{def:g1:plants-around-us:parts}{wortels} bijna niets, terwijl wind en zon het blijven
onttrekken --- het uitdroogprobleem van de pissebed, een hele winter
lang. De was en de naaldvorm van het \omterm{def:g6:life-through-seasons:trees}{groenblijvende} \omterm{def:g1:plants-around-us:parts}{blad} snijden die
verliezen terug tot wat de boom zich kan veroorloven. Een breed dun \omterm{def:g1:plants-around-us:parts}{blad}
laat zich niet zo goed afsluiten; het goedkopere antwoord van de
\omterm{def:g6:life-through-seasons:trees}{bladverliezende} boom is zijn dunne \omterm{def:g1:plants-around-us:parts}{bladeren} elke herfst af te schrijven
en de lente in plaats daarvan in \omterm{def:g6:life-through-seasons:buds}{knoppen} op te slaan.
\end{solution}

\begin{solution}{pb:g6:life-through-seasons:1}
\textbf{1.} De \omterm{def:g6:life-through-seasons:buds}{knoppen} op elke twijg --- winteropslag met de scheuten van
de volgende lente erin, de zomer ervoor gebouwd.

\textbf{2.} Tussen het maart- en het aprilbeeld (met de lengende dagen en
mildere ochtenden uit de aantekeningen). De boom antwoordt op de
stijgende \omterm{def:g6:exploring-our-environment:conditions}{temperatuur} en het lengende licht.

\textbf{3.} Bestuivende \omterm{prop:g3:sorting-living-things:groups}{insecten} --- vooral bijen --- moeten de \omterm{def:g3:flowers-fruits-seeds:parts}{bloemen}
bezoeken, zodat \omterm{def:g3:flowers-fruits-seeds:parts}{stuifmeel} de \omterm{def:g3:flowers-fruits-seeds:parts}{stampers} bereikt en de \omterm{def:g5:flower-to-fruit:steps}{bevruchting} volgt.
Een kastanje is het \omterm{def:g2:from-seed-to-plant:seed}{zaad} dat uit een bevruchte zaadknop gevormd is, in
zijn stekelige vruchtomhulsel.

\textbf{4.} \omterm{def:g6:life-through-seasons:trees}{Bladverliezend}. Zijn volgende lente rijdt de winterbeelden
door als \omterm{def:g6:life-through-seasons:buds}{knoppen} op de kale twijgen.

\textbf{5.} Onder het gazon, als bollen --- de ondergrondse opslag van de
\omterm{prop:g6:life-through-seasons:herbs}{overblijvende plant}.

\textbf{6.} Als \omterm{prop:g3:flowers-fruits-seeds:transform}{zaden} in de bodem: de strategie van de eenjarige --- de
\omterm{def:g1:plants-around-us:plant}{planten} zelf zijn in de herfst doodgegaan.

\textbf{7.} De zwaluw: in Afrika, als zichzelf getrokken, achter de
vliegende \omterm{prop:g3:sorting-living-things:groups}{insecten} aan. De bijen: in de kast, als wintertros, warm om
hun koningin heen.

\textbf{8.} Actief blijven. De ekster moet de magerste maanden door
voedsel zien te vinden --- precies het probleem waar de zwaluw met
vertrekken aan ontsnapt.

\textbf{9.} Bijvoorbeeld: de \omterm{def:g6:life-through-seasons:buds}{knoppen} springen open (maart--april) als de
dagen $11$--$13$ uur duren en de ochtenden boven de
\qty{8}{\degreeCelsius} komen; de \omterm{def:g1:plants-around-us:parts}{bladeren} verkleuren en vallen
(oktober--november) als de dagen korten en de ochtenden kouder worden.
Gebeurtenissen in de boom horen beeld na beeld bij veranderingen in de
\omterm{def:g6:exploring-our-environment:conditions}{omstandigheden}.

\textbf{10.} Het veldje is niet doodgegaan: het heeft zich in de opslag
teruggetrokken --- levend als bollen onder dezelfde plek. ``Weg'' verdient
altijd: weg \emph{waarheen}, als \emph{wat}?

\textbf{11.} Kale takken met \omterm{def:g6:life-through-seasons:buds}{knoppen} eraan; waarschijnlijk ergens een
ekster; eronder de bladlaag en kale narcissengrond met bollen erin
verstopt; klaproospitten in de aarde van het pad. De zekere weddenschappen
zijn de \omterm{def:g6:life-through-seasons:buds}{knoppen} en de kaalheid --- die volgen uit de \omterm{def:g6:life-through-seasons:trees}{bladverliezende}
machinerie van de boom en uit het bewijs van één jaar; de ekster is
waarschijnlijk, niet beloofd.

\textbf{12.} Het laat zien dat warmte de ontbrekende omstandigheid was
--- het licht is binnen in februari niet langer dan buiten, maar het
water en de kamer zijn warm, en dat volstond. Eerlijke proef: twee
twijgen van dezelfde boom, allebei in water bij hetzelfde raam, de een
in een warme kamer, de ander in een koude gang --- één verschil, de
warmte; als alleen de warme twijg uitloopt, staat het vonnis vast.
\end{solution}
